\chapter{Estructura de Retículo (Orden)}

\section{Relaciones de Orden}
El álgebra de Boole también puede formularse íntegramente prescindiendo de los operadores algebraicos, para fundamentarse de manera topológica o relacional mediante la Teoría del Orden. 

Para ello, introducimos una relación binaria $\le$ sobre el conjunto $\mathbb{B}$. 

\begin{definicion}{Conjunto Parcialmente Ordenado (Poset)}{poset}
Un conjunto $P$ equipado con una relación binaria $\le$ es un conjunto parcialmente ordenado (poset) si la relación satisface los siguientes axiomas para todo $a, b, c \in P$:
\begin{enumerate}
    \item \textbf{Reflexividad}: $a \le a$.
    \item \textbf{Antisimetría}: Si $a \le b$ y $b \le a$, entonces $a = b$.
    \item \textbf{Transitividad}: Si $a \le b$ y $b \le c$, entonces $a \le c$.
\end{enumerate}
\end{definicion}

En este contexto, escribiremos $a \ge b$ como sinónimo estricto de $b \le a$, y $a < b$ si $a \le b$ pero $a \ne b$.

\section{Ínfimo y Supremo}

Dado un poset $(P, \le)$, consideremos un par de elementos $a, b \in P$. 
\begin{itemize}
    \item Un elemento $u \in P$ es una \textbf{cota superior} de $\{a, b\}$ si $a \le u$ y $b \le u$. El \textbf{supremo} de $a$ y $b$, denotado como $a \sqcup b$, es la menor de todas sus cotas superiores (si existe).
    \item Un elemento $l \in P$ es una \textbf{cota inferior} de $\{a, b\}$ si $l \le a$ y $l \le b$. El \textbf{ínfimo} de $a$ y $b$, denotado como $a \sqcap b$, es la mayor de todas sus cotas inferiores (si existe).
\end{itemize}

\begin{definicion}{Retículo (Lattice)}{reticulo}
Un \textbf{retículo} es un conjunto parcialmente ordenado en el cual todo par de elementos tiene un supremo ($a \sqcup b$) y un ínfimo ($a \sqcap b$) definidos y únicos dentro del conjunto.
\end{definicion}

\section{Funtores de Equivalencia Estructural}

Al igual que ocurrió con los Anillos Booleanos, existe un isomorfismo total entre las Álgebras de Boole algebraicas y ciertos retículos con propiedades especiales.

\subsection{De Álgebra a Retículo}
Toda álgebra de Boole $(\mathbb{B}, \vee, \wedge)$ induce de forma natural un retículo definiendo la relación de orden parcial de la siguiente manera:
\[ a \le b \iff a \wedge b = a \]
Por las propiedades de absorción demostradas en capítulos previos, esta definición es lógicamente equivalente a su forma dual:
\[ a \le b \iff a \vee b = b \]
Bajo esta métrica de orden impuesta por las operaciones lógicas, es directo demostrar que el operador lógico $\vee$ computa exactamente el supremo topológico de los dos elementos ($a \sqcup b = a \vee b$), y que el operador lógico $\wedge$ computa su ínfimo ($a \sqcap b = a \wedge b$). A partir de este punto, el círculo se cierra y se revela que las operaciones algebraicas abstractas de Boole no son más que el cálculo de cotas topológicas en un espacio ordenado.

\subsection{De Retículo a Álgebra (Retículos Booleanos)}
No todo retículo es un álgebra de Boole. Para recuperar un álgebra de Boole plena desde la topología de orden puro, el retículo subyacente debe poseer tres propiedades restrictivas adicionales, dando lugar a lo que matemáticamente se conoce como \textbf{Retículo Booleano}:

\begin{teorema}{Axiomas del Retículo Booleano}{axiomas_reticulo_booleano}
Un retículo $(L, \le)$ es isomórfico a un Álgebra de Boole si y solo si es:
\begin{enumerate}
    \item \textbf{Acotado}: Existen elementos universales mínimo ($\bot$) y máximo ($\top$) tales que $\forall x \in L$, $\bot \le x \le \top$.
    \item \textbf{Distributivo}: El cálculo del ínfimo se distribuye sobre el cálculo del supremo, y viceversa.
    \item \textbf{Complementado}: Para cada elemento $a \in L$ existe un elemento único $b \in L$ (denominado su complemento, $\neg a$) tal que el ínfimo de ambos es $\bot$ y su supremo es $\top$.
\end{enumerate}
\end{teorema}

\section{Implicación para las Álgebras Finitas}
Esta profunda visión topológica como un retículo complementado y distributivo es la base fundamental que nos permitirá visualizar las álgebras de Boole en el espacio mediante \textit{Diagramas de Hasse}. 

Más importante aún, esta estructura de orden estricto impone un límite severo a la cardinalidad. Si un álgebra de Boole tiene un número finito de elementos, su retículo subyacente se construirá forzosamente como el conjunto potencia de sus átomos (los elementos inmediatamente superiores al $\bot$). Esto revela, por el Teorema de Representación de Stone para retículos finitos, que \textbf{toda álgebra de Boole finita debe tener forzosamente $2^n$ elementos}. 

Este hecho es el puente de entrada y la justificación absoluta para adentrarnos en el siguiente capítulo: el estudio y la forma de las Álgebras de Boole Finitas.
